\begin{definition} \label{definition:geometric_point_of_a_scheme}
Let $S$ be a \CrefAndHyperrefIfExist{definition:scheme}{scheme}. There are two inequivalent conventions for the definition of geometric points of $S$:
\begin{enumerate}
    \item A \hldef{(algebraic) geometric point of $S$} is a morphism $\overline{s} : \Spec(\Omega) \to S$ where $\Omega$ is a \CrefAndHyperrefIfExist{definition:separably_closed_field}{separably closed field}. 

    \item A \hldef{(separable) geometric point of $S$} is a morphism $\overline{s} : \Spec(\Omega) \to S$ where $\Omega$ is an \CrefAndHyperrefIfExist{definition:algebraically_closed_field}{algebraically closed field}. 
\end{enumerate}
Care should be taken to specify which convention is used in a given context. In general, one should use the algebraic notion of geometric points for algebraic discussions, say those  involving the Nullstellensatz, classical intersection theory/degrees, ``\CrefAndHyperrefIfExist{definition:geometric_property_of_a_scheme_over_a_field}{geometric properties}'' of schemes over fields, etc., and the separable notion of geometric points for arithmetic discussions, say those involving Galois theory, Galois modules and cohomology, \'etale sites, \'etale fundamental groups, stalks of \'etale sheaves, etc. 

The two theory of the two conventions are usually equivalent when ``purely inseparable phenomena'' are not present and usually inequivalent when purely inseparable phenomena are present. 
\end{definition}